\chapter{Contexte général du projet et analyse des besoins}

\section{Introduction}

Ce chapitre ouvre le rapport en présentant le cadre général du projet et en précisant les besoins auxquels répond la solution développée pendant le stage. Après l'introduction générale, qui a introduit les grands thèmes liés à la digitalisation des ressources humaines et à l'aide à la décision, il s'agit ici d'entrer dans le concret : le contexte de l'entreprise d'accueil, la problématique rencontrée, les objectifs visés et les fonctionnalités attendues de la plateforme. L'objectif est de donner au lecteur une vision claire avant d'aborder, dans les chapitres suivants, l'étude bibliographique, le cahier des charges, la conception, la réalisation et les tests.

Le projet porte sur \textbf{Strategic Copilot}, une plateforme web d'aide à la décision pour le pilotage des projets et des talents. Elle s'inscrit dans l'offre plus large \textbf{Augmented Talents} et s'adresse à trois types d'utilisateurs : les \textbf{managers}, qui suivent leur portefeuille de projets et leur équipe ; les \textbf{responsables RH}, qui gèrent les compétences, les formations et les demandes des managers ; et les \textbf{talents}, qui consultent leurs missions, leurs tâches et leur charge de travail. J'ai conçu et développé l'application \texttt{copilot-ui} de manière à regrouper ces trois espaces dans une même solution, avec des droits d'accès adaptés à chaque profil. Le développement a débuté par l'\textbf{espace Manager}, aujourd'hui le plus abouti en termes de fonctionnalités, avant d'étendre progressivement les modules RH et Talent ; l'ensemble de la plateforme reste néanmoins couvert par le même socle technique.

Pour rédiger cette analyse, je me suis appuyée sur le document \emph{PFE Agentic AI Strategic Talents}, sur le cahier des charges fourni en entreprise, sur les échanges avec les encadrants, et sur l'étude du code source du dépôt \texttt{frontend-ui}. Seules les fonctionnalités réellement présentes ou clairement prévues dans ces documents y sont mentionnées. Ce chapitre pose ainsi les bases nécessaires pour comprendre la suite du rapport et pour relier chaque besoin métier aux modules concrets de la plateforme.

\section{Présentation de l'entreprise d'accueil}

Le stage de fin d'études s'est déroulé au sein de \textbf{AphelionX Innovations}, entreprise marocaine basée à Tétouan, spécialisée dans le développement logiciel, le conseil en systèmes d'information, l'automatisation des processus métier fondée sur l'intelligence artificielle et l'accompagnement des organisations dans leur transformation digitale. Ce positionnement place l'entreprise à l'intersection de l'ingénierie applicative et de l'innovation par les agents intelligents, ce qui constitue un cadre particulièrement adapté à un projet tel que Strategic Copilot, où la valeur repose autant sur la qualité de l'interface que sur la fiabilité des workflows backend et des modèles d'analyse.

Dans ce contexte, AphelionX développe la solution \textbf{Srategique Copilot}, plateforme modulaire orientée vers le pilotage conjoint des projets et des compétences. Strategic Copilot en est la brique d'intelligence décisionnelle : elle ne se substitue pas aux outils de gestion opérationnelle classiques, mais les prolonge en consolidant des signaux projet et RH, en détectant des risques et en proposant des recommandations que les décideurs peuvent examiner, accepter ou rejeter. L'infrastructure technique que j'ai mise en place repose sur une instance n8n dédiée (\texttt{n8nprod.aphelionxinnovations.com}), qui expose l'ensemble des webhooks métier et orchestre les workflows agentiques, ainsi que sur une base de données PostgreSQL hébergée via \textbf{Supabase} pour la persistance relationnelle et, selon les cas, le stockage de fichiers (rapports PDF, avatars). Le front-end React consomme exclusivement ces services via des appels HTTP sécurisés ; il n'effectue pas de calculs critiques de scores de viabilité, de risques ou de recommandations en local, conformément au principe API-first rappelé dans le typage métier (\texttt{types/crud-domain.ts}).

L'encadrement en entreprise a été assuré par Monsieur \textbf{Hicham Taoufik}, Monsieur \textbf{Youssef Moufak}, Monsieur \textbf{Hamza Sadouk} et Monsieur \textbf{Mouhy Eddine Saadani}, qui ont orienté les choix fonctionnels, validé les priorités de livraison et accompagné l'intégration des composants Copilot et des agents métier documentés dans le cahier Agentic AI. L'encadrement pédagogique de Madame \textbf{Soundouss Abroun}, au sein de l'École Marocaine des Sciences de l'Ingénieur (EMSI) de Tanger, a garanti la rigueur méthodologique, la structuration du rapport et l'alignement du livrable sur les exigences académiques du diplôme d'Ingénieur d'État en Informatique et Réseaux, option MIAGE. Le projet s'inscrit ainsi à l'intersection d'une démarche industrielle de type produit SaaS et d'une démarche d'ingénierie logicielle structurée, où la traçabilité des décisions et la clarté des interfaces sont considérées comme des exigences au même titre que la performance technique.

\section{Contexte du stage}

Le stage s'est déroulé du \textbf{2 février 2026 au 2 juillet 2026}, soit une durée de \textbf{cinq mois}, dans le cadre du Projet de Fin d'Études de l'EMSI Tanger en vue de l'obtention du diplôme d'Ingénieur d'État en Informatique et Réseaux, option MIAGE. J'ai assuré personnellement l'ensemble des grandes étapes du projet Strategic Copilot : analyse des besoins à partir du document \emph{PFE Agentic AI Strategic Talents}, conception de l'architecture fonctionnelle et technique, développement de l'interface avec \textbf{React}, \textbf{Vite} et \textbf{TypeScript}, mise en place des workflows backend avec \textbf{n8n}, modélisation et alimentation de la base \textbf{Supabase PostgreSQL}, intégration des modèles \textbf{Llama} pour les fonctions d'intelligence artificielle, puis tests et validation sur l'environnement de préproduction.

Sur le plan académique, ce stage représente la synthèse des compétences acquises en conception, en développement web, en intégration de services et en rédaction technique. Sur le plan professionnel, il répond au besoin d'AphelionX de disposer d'une plateforme unifiée, multilingue et sécurisée, capable de restituer les analyses produites par les agents métier sans les déformer. J'ai développé les modules des trois espaces --- \textbf{Manager}, \textbf{RH} et \textbf{Talent} --- au sein d'une même application routée dans \texttt{main.tsx}. L'espace Manager (tableau de bord, portefeuille projet, Mission Control, équipe, risques, demandes RH, rapports, notifications, Copilot et Helper) a été structuré en premier et reste aujourd'hui le plus riche fonctionnellement ; les espaces RH et Talent disposent pour leur part de parcours complets (tableaux de bord, référentiels, demandes, missions, charge, compétences, formations), intégrés aux mêmes services d'authentification et aux mêmes workflows n8n.

Le déroulement du projet a suivi une logique incrémentale. Chaque livraison a été validée par des tests d'intégration contre les webhooks n8n, par des ajustements des contrats TypeScript (\texttt{api.types.ts}) et par des retours des encadrants sur l'ergonomie des écrans. L'environnement local repose sur \textbf{Vite} avec un proxy vers n8n pour contourner les contraintes CORS. Les rôles, les statuts utilisateurs et l'authentification JWT sont entièrement gérés côté serveur ; le front-end n'ouvre pas de parcours d'auto-inscription, le provisionnement des comptes relevant du périmètre RH.

\section{Problématique}

La problématique métier, formulée dans le document \emph{PFE Agentic AI Strategic Talents} et confirmée par les retours recueillis durant le stage, peut être analysée comme une tension structurelle entre la \textbf{richesse des données} disponibles en entreprise et la \textbf{pauvreté de leur mise en scène décisionnelle}. Dans de nombreuses organisations, les informations relatives aux projets (jalons, budgets, affectations), aux talents (compétences, contrats, charge), aux demandes RH (formation, mobilité, recrutement, réallocation) et aux signaux de risque (surcharge, retards, écarts de compétences) sont produites par des outils hétérogènes : feuilles de calcul, messageries, tickets, systèmes RH partiellement connectés. Cette fragmentation limite la visibilité globale, retarde la détection des dérives et rend la prise de décision plus coûteuse en temps et en coordination.

Pour le \textbf{manager}, la difficulté consiste à piloter simultanément un portefeuille de projets et une équipe de collaborateurs sans disposer d'une vue consolidée actualisée. Les retards sur un livrable, une montée en charge imprévue sur un profil clé ou l'approche d'une fin de contrat peuvent ne être identifiés que lorsque les marges de manœuvre sont réduites. Pour le \textbf{responsable RH}, la problématique est différente mais connexe : il doit maintenir un référentiel de compétences fiable, détecter les écarts critiques (\emph{skill gaps}), instruire des plans de formation et de mobilité, tout en traitant des demandes émises par les managers avec des types métier variés (\texttt{skill\_gap}, \texttt{reallocation}, \texttt{training}, \texttt{overload}, \texttt{recruitment}) sans toujours disposer d'une vision temps réel de l'impact sur les projets en cours. Pour le \textbf{talent}, l'enjeu est la transparence : comprendre ses missions, ses tâches, sa charge de travail et ses axes de développement, sans accéder à des données hors périmètre ni subir une opacité sur les décisions qui le concernent.

À ces difficultés s'ajoute l'émergence de l'\textbf{intelligence artificielle générative} comme levier d'analyse et de synthèse. Les modèles de langage, lorsqu'ils sont correctement encapsulés dans des workflows contrôlés (orchestration n8n, paramètre \texttt{use\_ai}, modèles tels que \textbf{Llama} mobilisés côté serveur selon le cadrage du projet), peuvent accélérer la détection de signaux faibles et la formulation de recommandations. Toutefois, leur adoption en contexte managérial et RH exige une \textbf{explicabilité} des sorties, une séparation claire entre corrélation et décision, et une traçabilité des arbitrages humains. Strategic Copilot ne prétend pas remplacer le jugement des décideurs ; il vise une aide à la décision \emph{augmentée}, en transformant des analyses multi-agents complexes en indicateurs, alertes, scores de viabilité et propositions d'action que l'utilisateur peut valider, modifier ou rejeter.

La problématique du stage peut donc se résumer ainsi : comment concevoir et livrer une interface web unifiée, sécurisée par rôle, qui centralise les données RH et projet, expose les résultats des agents d'analyse et de risque, facilite la collaboration manager--RH--talent, et respecte une architecture où le calcul métier critique reste côté serveur ? Cette question oriente l'ensemble des objectifs et des besoins formalisés dans les sections suivantes.

\section{Objectifs du projet}

Les objectifs du projet de fin d'études découlent directement de la problématique et du cahier des charges Agentic AI. Ils se déclinent en quatre dimensions complémentaires : métier, technique, agentique et utilisateur.

Sur le plan \textbf{métier}, il s'agit de \textbf{centraliser} les données relatives aux projets et aux talents dans une expérience cohérente, de \textbf{détecter} proactivement les risques liés aux surcharges, aux retards, aux écarts de compétences et aux échéances contractuelles, de \textbf{fournir} des recommandations issues de l'intelligence artificielle de manière structurée et interprétable, et de \textbf{faciliter} la collaboration entre managers, responsables RH et talents via des flux de demandes, d'alertes et de notifications traçables. La vision produit affichée dans le dépôt résume cette ambition : transformer des analyses complexes portant sur les projets, les compétences, les charges et les risques en décisions managériales claires, explicables et directement actionnables.

Sur le plan \textbf{technique}, l'objectif est de livrer une application \textbf{React 19} typée en \textbf{TypeScript}, construite avec \textbf{Vite 7}, stylée en \textbf{Tailwind CSS 4} et fondée sur des composants accessibles (\textbf{React Aria Components}, kit \textbf{Untitled UI}), consommant les webhooks \textbf{n8n} et la base \textbf{PostgreSQL} via Supabase sans logique décisionnelle critique dans le navigateur. L'authentification repose sur des jetons \textbf{JWT} (\texttt{/webhook/login}, \texttt{/webhook/refresh}) ; les routes applicatives sont protégées par rôle (\texttt{ProtectedRoute}). L'internationalisation (\textbf{i18next}, français, anglais, arabe) et la gestion d'état serveur (\textbf{TanStack React Query}) visent une expérience fluide malgré la latence potentielle des workflows longs.

Sur le plan \textbf{agentique}, l'objectif est d'exposer de manière homogène les résultats des agents documentés : \textbf{Project Analysis Agent} (\texttt{WF\_Project\_Analysis}, analyste projet), \textbf{Talent Matching Agent} (\texttt{WF\_Talent\_Matching}, appariement compétences--besoins), \textbf{Risk \& Weak Signal Agent} (\texttt{WF\_Project\_Risk}, surveillance et scan Watchdog), et \textbf{Strategic Orchestrator Agent} (\texttt{WF\_Project\_Viability}, viabilité et scénarios). Le Helper conversationnel (\texttt{POST /webhook/api/helper/chat}) et le Copilot global (\texttt{CopilotProvider}, \texttt{CopilotPanel}) constituent les points d'interaction naturels avec ces capacités.

Sur le plan \textbf{utilisateur}, l'objectif est de proposer trois espaces routés (\texttt{/workspace/manager}, \texttt{/workspace/rh}, \texttt{/workspace/talent}), chacun adapté aux décisions de son profil, avec des tableaux de bord analytiques, des parcours de traitement des alertes et des demandes, et des mécanismes de reporting (génération PDF, envoi par e-mail via \texttt{/webhook/reports/send-email} lorsque le workflow est disponible). Les critères de réussite incluent la disponibilité des parcours critiques testés en préproduction : connexion, pilotage projet, scan Watchdog, demande RH, simulation what-if, enregistrement d'une décision dans le journal Copilot, et cohérence des données affichées avec les réponses des workflows n8n.

\section{Présentation générale de Strategic Copilot}

\subsection{Vision et positionnement}

Strategic Copilot se définit comme une plateforme \textbf{SaaS intelligente} de pilotage stratégique des talents et des projets. Son positionnement au sein d'Augmented Talents est celui d'un \textbf{copilote décisionnel} : il ne remplace ni l'outil de gestion de projet classique ni le système d'information RH, mais les enrichit en agrégeant des indicateurs, en détectant des risques et en proposant des scénarios d'arbitrage. La vision, rappelée dans le README du dépôt, insiste sur la conversion d'analyses multi-agents en décisions actionnables (\emph{Transform complex multi-agent analysis into clear, actionable decisions}).

Le cœur métier du pilotage projet repose sur un \textbf{score de viabilité} et une \textbf{décision synthétique} produits côté orchestrateur n8n, selon une pondération documentée dans le cahier des charges : Skills Fit (40\,\%), Capacité (35\,\%), Budget (25\,\%). Le score est traduit en recommandation \texttt{Continue} (score $\geq 7{,}5$), \texttt{Adjust} (entre 4 et 7,5) ou \texttt{Stop} ($< 4$). L'interface se limite à afficher les champs \texttt{viability\_score}, \texttt{risk\_score} et \texttt{decision} renvoyés par le serveur, ce qui garantit une cohérence entre les utilisateurs et évite les divergences de calcul. Le positionnement « SaaS » implique par ailleurs une séparation nette entre la couche présentation (build Vite, déploiement possible sur Vercel avec relais des webhooks) et la couche orchestration--données (n8n, PostgreSQL), chaque couche pouvant évoluer indépendamment tant que les contrats d'API restent stables.

\subsection{Les trois espaces utilisateurs}

L'application structure l'expérience en trois espaces utilisateurs distincts mais articulés autour des mêmes entités métier (projets, talents, compétences, alertes, demandes). Le routage est centralisé dans \texttt{main.tsx} ; l'accès est filtré par rôle JWT via \texttt{ProtectedRoute} et des layouts dédiés qui assurent une navigation homogène au sein de chaque perspective.

\paragraph{Espace Manager.}
L'espace \texttt{/workspace/manager} concentre les besoins du pilotage opérationnel et stratégique. Le \textbf{tableau de bord} (\texttt{DashboardPage}) restitue des indicateurs consolidés et, lorsque le backend les fournit, des blocs issus des agents Matchmaker et Analyst. Le \textbf{portefeuille projets} (\texttt{ProjectsPage}) permet de parcourir l'ensemble des projets suivis et d'ouvrir le modal \textbf{Mission Control} (\texttt{ProjectMissionControlModal}), véritable centre de pilotage à 360 degrés d'un projet : synthèse, équipe, tâches, risques, simulation what-if et journal des décisions liées. La gestion d'\textbf{équipe} (\texttt{TeamPage}, \texttt{TalentDetailPage}) offre une lecture fine des profils, des allocations et des alertes, avec possibilité de déclencher un scan Watchdog (\texttt{POST /webhook/api/watchdog/scan}). La page \textbf{Risques \& alertes} (\texttt{RisksPage}) propose une vue transverse des signaux critiques, complétée par la page \textbf{Notifications} (\texttt{NotificationsPage}) qui consomme les workflows de type \texttt{WF\_Manager\_Notifications}. Les \textbf{demandes RH} (\texttt{ManagerRhRequestsPage}, composants \texttt{rh-requests}) et les \textbf{demandes talent} (\texttt{TalentRequestsPage}) matérialisent les flux collaboratifs vers les autres acteurs. Le \textbf{journal des décisions} (\texttt{DecisionLogPage}), les \textbf{rapports} (\texttt{ReportsPage}, génération board-pack, historique, partage par e-mail) et le \textbf{Helper} (\texttt{HelperChatPage}) complètent l'espace. Le \textbf{Copilot} contextuel (\texttt{CopilotProvider}, panneau global \texttt{CopilotPanel}) accompagne l'utilisateur selon le contexte consulté (tableau de bord, détail projet, staffing). Cet espace concentre aujourd'hui le plus grand nombre de fonctionnalités avancées ; il a constitué le point de départ du développement, sans exclure pour autant la réalisation des modules RH et Talent décrits ci-dessous.

\paragraph{Espace RH.}
L'espace \texttt{/workspace/rh} répond aux besoins de gouvernance des ressources humaines et de traitement des demandes managers. Le \textbf{tableau de bord RH} (\texttt{RhDashboardPage}) offre une vue agrégée des indicateurs RH. Le module \textbf{employés / talents} (\texttt{RhEmployeesPage}) permet de consulter et d'administrer le référentiel des collaborateurs. Le \textbf{catalogue de compétences} (\texttt{RhSkillsCatalogPage}) et la page des \textbf{écarts critiques} (\texttt{RhCriticalGapsPage}) soutiennent la politique de compétences. Les \textbf{plans de formation} (\texttt{RhTrainingPlansPage}) et la \textbf{mobilité interne} (\texttt{RhMobilityPage}) adressent le développement et la circulation des talents. Les \textbf{alertes organisationnelles} (\texttt{RhOrgAlertsPage}) complètent la surveillance transverse. Le traitement des \textbf{demandes managers} (\texttt{RhManagerRequestsPage}, \texttt{rh-manager-requests-page}) permet d'accepter, rejeter ou faire progresser les demandes typées. Des vues \textbf{comptes} et \textbf{sessions} (\texttt{RhAccountsWorkspacePage}, \texttt{RhSessionsWorkspacePage}) assurent l'administration des accès, tandis que les \textbf{rapports RH} (\texttt{RhReportsWorkspacePage}) et une vue \textbf{projets} (\texttt{ProjectsPage} dans le layout RH) maintiennent l'alignement avec le portefeuille manager. J'ai également développé cet espace : tableaux de bord, pages métiers et appels API dédiés (\texttt{/workspace/rh/*}) sont opérationnels et connectés aux workflows n8n correspondants.

\paragraph{Espace Talent.}
L'espace \texttt{/workspace/talent} adopte le point de vue du collaborateur. Le \textbf{tableau de bord personnel} (\texttt{TalentDashboardPage}) synthétise la situation individuelle. Les pages \textbf{Mes projets} (\texttt{TalentProjectsPage}), \textbf{Mes tâches} (\texttt{TalentTasksPage}) et \textbf{Charge de travail} (\texttt{TalentWorkloadPage}) rendent visibles les affectations et la charge associée. La page \textbf{Compétences} (\texttt{TalentSkillsPage}) et celle des \textbf{Formations} (\texttt{TalentTrainingPage}) soutiennent le développement professionnel. Les \textbf{notifications} (\texttt{TalentNotificationsPage}) et le \textbf{profil} (\texttt{TalentProfileWorkspacePage}) ferment la boucle en offrant une lecture individuelle des informations qui, en amont, alimentent les analyses vues par le manager et le RH. Les données talent sont chargées via des endpoints agrégés de type \texttt{/api/workspace/talent/*}, ce qui garantit un périmètre strict au collaborateur connecté.

La coexistence de ces trois espaces dans une même application n'est pas une simple juxtaposition d'écrans : elle traduit un modèle organisationnel où les décisions se prennent en chaîne. Le manager identifie un risque ou un besoin, formalise une demande RH, le responsable RH instruit et met à jour le référentiel, le talent constate l'impact sur ses missions et sa charge. Strategic Copilot se veut le médiateur numérique de cette chaîne.

\subsection{Fonctionnalités transverses}

Plusieurs capacités traversent les trois espaces et constituent l'identité produit de Strategic Copilot au-delà des menus spécifiques à chaque rôle.

L'\textbf{authentification sécurisée par rôle} (\texttt{AuthProvider}, API \texttt{auth.api.ts}) conditionne l'accès aux workspaces et évite les fuites de données entre profils. Les \textbf{tableaux de bord analytiques}, appuyés sur Recharts pour certaines visualisations manager, traduisent des agrégats serveur en graphiques et cartes KPI lisibles. La \textbf{détection automatique des risques} et le mécanisme d'\textbf{alertes et notifications} s'appuient sur les workflows Watchdog et Manager Notifications ; l'interface permet l'acquittement, la résolution ou l'ignorance selon que l'alerte est liée à une entrée \texttt{risk\_alerts} ou à une notification seule.

Les \textbf{recommandations IA} sont restituées via le Copilot, le panneau Matchmaker, les synthèses Analyst et les réponses du Strategist lorsque \texttt{use\_ai} est activé dans les appels n8n. Le \textbf{simulateur what-if} (\texttt{WhatIfProvider}, \texttt{POST /webhook/api/project/what-if}) permet de tester l'impact de modifications d'allocation, de délai ou de compétences avant arbitrage. La \textbf{génération de rapports PDF} et l'\textbf{envoi par e-mail} (\texttt{use-reports-n8n}, \texttt{sendReportEmail}) supportent la gouvernance et le partage formel des analyses. Le \textbf{journal des décisions} (\texttt{DecisionLogPage}, API \texttt{/webhook/manager/copilot-decisions}) assure la traçabilité des choix Continue / Adjust / Stop et des motifs associés. Enfin, le \textbf{Helper conversationnel} propose un dialogue contextualisé avec suggestions d'actions (\texttt{assign}, \texttt{training}, \texttt{review}, \texttt{risk}), complétant le Copilot plus orienté synthèse et scores.

Ces fonctions transverses incarnent l'architecture multi-agents : le front-end orchestre les appels et la restitution, tandis que n8n enchaîne lecture PostgreSQL, règles métier et inférences Llama sans exposer les secrets de modèle au client.

\section{Analyse des besoins}

\subsection{Acteurs et cas d'usage}

L'analyse des besoins commence par l'identification des acteurs et de leurs intentions principales, telles qu'elles ressortent du cahier Agentic AI et des parcours implémentés.

Le \textbf{manager} souhaite connaître l'état de son portefeuille projet, identifier les projets à risque, affecter ou réaffecter des talents, formuler des demandes RH ou talent, simuler des scénarios avant engagement, consigner ses décisions et consulter des rapports de synthèse. Ses cas d'usage prioritaires incluent l'ouverture du Mission Control sur un projet critique, le déclenchement d'un scan Watchdog sur un collaborateur en surcharge, la création d'une demande de type \texttt{training} ou \texttt{reallocation}, et l'exécution d'une simulation what-if avant de valider un arbitrage Strategist.

Le \textbf{responsable RH} souhaite maintenir le référentiel talents et compétences, détecter les écarts critiques, planifier formations et mobilités, traiter les demandes managers dans des délais maîtrisés, surveiller les alertes organisationnelles et administrer les comptes utilisateurs. Ses cas d'usage couvrent l'instruction d'une demande manager (acceptation, rejet, mise en progression), la consultation du catalogue de compétences, l'analyse des gaps sur un périmètre métier et la génération de rapports RH.

Le \textbf{talent} souhaite consulter ses missions, mettre à jour l'état de ses tâches lorsque le workflow le permet, suivre sa charge, visualiser ses compétences et formations, recevoir des notifications pertinentes et gérer son profil. Le système doit garantir qu'il n'accède pas aux données d'autres collaborateurs ni aux vues stratégiques réservées au manager.

Les processus transverses documentés (P0 à P5 dans le cahier Agentic AI) structurent ces cas d'usage : synchronisation des données (\texttt{WF\_Data\_Sync}), analyse projet, matching talents, analyse des risques, orchestration stratégique et simulation. Chaque processus se traduit par un ou plusieurs webhooks HTTP déclenchés depuis l'interface ou par des tâches planifiées côté n8n.

\subsection{Besoins fonctionnels}

Les besoins fonctionnels expriment les services attendus du système. Ils ont été validés par confrontation entre le cahier des charges et les modules présents dans \texttt{copilot-ui}.

Le système doit assurer l'\textbf{authentification et la gestion de session} (connexion, renouvellement, déconnexion, profil, changement de mot de passe). Il doit permettre le \textbf{pilotage projet} : liste, détail versionné (\texttt{wmp-detail-v1}), affectations (\texttt{wmp-assign-v1}, \texttt{wmp-unassign-v1}), tâches (\texttt{wmt-*-v1}), représentation du cycle de vie (\texttt{ProjectLifecycleStepper}). Il doit couvrir la \textbf{gestion des talents et compétences} côté RH et manager : référentiel, fiche talent, catalogue \texttt{Skill}, liaisons \texttt{TalentSkill}. Il doit exposer les \textbf{risques et alertes} : liste par projet, notifications manager, actions de résolution ou d'acquittement. Il doit délivrer les \textbf{recommandations IA} via Copilot, Matchmaker, Analyst, Helper et Strategist. Il doit gérer les \textbf{demandes RH} (création manager, liste, traitement RH par PATCH). Il doit proposer la \textbf{simulation what-if} et les \textbf{rapports} (génération, historique, envoi e-mail). Il doit enfin offrir les \textbf{espaces talent} (dashboard, projets, tâches, charge, compétences, formations, notifications, profil).

Le tableau~\ref{tab:besoins-ch1} synthétise ces besoins par domaine et indique la trace dans l'application livrée.

\begin{table}[htbp]
\centering
\caption{Besoins fonctionnels et trace dans l'application}
\label{tab:besoins-ch1}
\begin{tabular}{@{}p{2.6cm}p{5.4cm}p{4.6cm}@{}}
\toprule
\textbf{Domaine} & \textbf{Besoin} & \textbf{Réalisation (copilot-ui)} \\
\midrule
Authentification & Connexion sécurisée par rôle & \texttt{AuthProvider}, \texttt{/webhook/login} \\
Portefeuille & Vue consolidée projets & \texttt{DashboardPage}, dashboard manager \\
Projet & Pilotage détaillé 360° & \texttt{ProjectMissionControlModal} \\
Équipe & Suivi talents, scan risques & \texttt{TeamPage}, Watchdog \\
RH & Référentiel, écarts, demandes & Routes \texttt{/workspace/rh/*} \\
Talent & Missions, charge, compétences & Routes \texttt{/workspace/talent/*} \\
IA & Synthèse, chat, arbitrage & \texttt{CopilotProvider}, Helper, Strategist \\
Simulation & Scénarios what-if & \texttt{WhatIfProvider} \\
Reporting & PDF, e-mail, historique & \texttt{ReportsPage}, \texttt{reports.api} \\
Notifications & Alertes, acquittement & \texttt{NotificationsPage} \\
\bottomrule
\end{tabular}
\end{table}

\subsection{Besoins non fonctionnels}

Les besoins non fonctionnels encadrent la qualité de la solution et conditionnent son adoption en entreprise.

En matière de \textbf{performance}, l'application doit tolérer la latence des workflows n8n grâce à des délais configurables (\texttt{VITE\_HTTP\_TIMEOUT\_MS}, \texttt{VITE\_PORTFOLIO\_TIMEOUT\_MS}) et à TanStack React Query (cache, invalidation ciblée après mutations, par exemple \texttt{invalidateManagerRiskQueries}). Les états de chargement, d'erreur et de vide doivent être explicites pour éviter toute ambiguïté sur l'état du système.

En matière de \textbf{sécurité}, les jetons JWT ne doivent pas exposer les clés des modèles IA ; les routes sont filtrées par rôle ; le client gère le renouvellement de session sur réponse 401. Les données talent ne sont accessibles qu'au talent connecté via les endpoints workspace dédiés.

En matière de \textbf{fiabilité}, les parseurs tolèrent les réponses n8n enveloppées (\texttt{json}, \texttt{body}, tableaux \texttt{notifications} ou \texttt{items}) ; les messages d'erreur sont orientés utilisateur (\texttt{readUserFacingApiErrorMessage}).

En matière de \textbf{maintenabilité}, TypeScript structure les contrats dans \texttt{api.types.ts} ; les chemins d'API sont externalisés en variables \texttt{VITE\_*} pour accompagner l'évolution des webhooks versionnés (\texttt{wmp-detail-v1}, etc.).

En matière d'\textbf{utilisabilité}, l'interface est responsive, accessible (React Aria), traduite (i18next) et cohérente visuellement (Tailwind CSS 4, composants Untitled UI). En matière d'\textbf{explicabilité IA}, les scores et décisions affichés proviennent du serveur ; l'interface rappelle, sur certains panneaux, que les calculs ne sont pas refaits côté client.

En matière de \textbf{traçabilité}, le journal des décisions, l'historique des rapports et le suivi des demandes RH permettent un audit a posteriori des choix organisationnels.

\subsection{Contraintes identifiées}

Plusieurs contraintes ont guidé les choix de conception et expliquent certains arbitrages du stage.

La contrainte de \textbf{temps limité} (cinq mois) a conduit à organiser le développement en commençant par l'espace Manager et les parcours les plus critiques pour la démonstration, tout en livrant les espaces RH et Talent sur la même architecture applicative.

La \textbf{complexité des processus RH} (typologie des demandes, SLA, validation multi-niveaux) exige des workflows n8n maintenables et des contrats d'API stables ; toute évolution métier doit être répercutée côté serveur avant l'interface.

La contrainte d'\textbf{explicabilité des recommandations IA} interdit de présenter des sorties de modèle comme des vérités définitives ; l'interface doit laisser la responsabilité au décideur et journaliser les arbitrages.

La \textbf{performance des workflows n8n} (analyses longues, appels Llama) impose des timeouts élevés et des indicateurs de progression ; le front-end ne doit pas bloquer l'expérience sans retour visuel.

La contrainte de \textbf{cohérence et qualité des données} rappelle que la valeur de Strategic Copilot dépend de l'exhaustivité des enregistrements PostgreSQL (affectations, compétences, dates de contrat) ; une alerte sans \texttt{risk\_alert\_id} correctement renseigné ne peut pas être résolue via le PATCH \texttt{risk-alerts} et doit être traitée par acquittement notification, comportement que l'interface gère désormais explicitement.

Enfin, les contraintes architecturales \textbf{API-first} et \textbf{pas de calcul front} (\texttt{crud-domain.ts}) imposent que toute nouvelle règle de viabilité, de matching ou de risque transite par n8n avant affichage. Ces contraintes, loin de freiner le projet, ont permis de maintenir une séparation claire des responsabilités et de préparer les évolutions backend sans remise en cause systématique des écrans.

\section{Conclusion}

Ce chapitre a présenté le contexte général du projet Strategic Copilot et l'analyse des besoins qui ont guidé mon travail de fin d'études. L'entreprise d'accueil AphelionX Innovations, le cadre du stage à l'EMSI Tanger et l'encadrement professionnel et pédagogique ont été rappelés. La problématique de fragmentation des données, de détection tardive des dérives et d'arbitrages peu traçables a été reliée aux objectifs du livrable.

La présentation de la plateforme a montré comment Strategic Copilot réunit trois espaces complémentaires --- \textbf{Manager}, \textbf{RH} et \textbf{Talent} --- autour d'une architecture multi-agents (analyse projet, matching talents, risques, orchestrateur stratégique), d'une stack React--n8n--PostgreSQL--Supabase et d'une intelligence artificielle orchestrée côté serveur (modèles Llama). J'ai personnellement porté l'ensemble du cycle projet, de l'analyse des besoins à la validation, en commençant par l'espace Manager, le plus avancé à ce stade, sans que mon travail se limite à ce seul périmètre.

L'analyse des besoins a formalisé les attentes des acteurs, les besoins fonctionnels et non fonctionnels, ainsi que les contraintes observées dans le code. Les chapitres suivants détailleront l'étude bibliographique, le cahier des charges, la conception, la réalisation et les tests de la solution que j'ai développée.
